% 019 两客户端同一业务：在现行 UC-01～15 后续编，不沿用旧版未引用的 UC-16～24 草稿。
\subsection*{3.2 补充：网页与微信小程序共用顾客业务}
\noindent 微信小程序是顾客的第二个边界界面，不是第二套库存或第二个顾客账号。网页顾客与已显式绑定的小程序顾客映射到同一内部顾客 ID；两端读取同一购物车、预约、单件库存与责任进度。门店专员、经理和技师继续在网页端完成实体交接及鉴定。微信平台只提供临时身份凭证与正式小程序码，不拥有 BSR 订单裁决权。
\BSRDiagram{019-usecases}{U-MINI 两客户端协作用例：顾客、微信平台与门店专员均在系统边界外；所有用例为椭圆}{fig:u-mini}

\Needspace{11\baselineskip}
\subsubsection{3.2.16 UC-16 查看网页扫码入口}\label{uc:16}
\begin{BSRUsecase}{UC-16 查看网页扫码入口}
主执行者 & 顾客。\\
辅执行者 & 微信平台提供正式小程序码；网站访客无需账号也可看入口状态。\\
涉众利益 & 顾客知道是否能从微信进入同一租赁服务；运营者不向顾客展示无效码或错误 AppID。\\
触发与前置 & 顾客打开网页入口；服务端已有明确的入口启停状态。\\
成功后置 & 仅在真实码已按配置核验且入口就绪时展示扫码图片；码不传订单、个人资料或网页会话。\\
失败后置 & 配置缺失、失效或停用时显示“尚未开放”及网页继续操作入口，不展示开发码或伪码。\\
基本流程 & 1 网页向服务端查询入口状态。2 系统验证环境、发行状态、码文件及允许的场景。3 就绪时返回受控图片与公开说明。4 顾客使用微信扫码打开小程序。5 小程序独立执行 UC-17 身份验证。\\
扩展流程 & 2a 无正式 AppID、未审核、码失效或校验失败：入口关闭。4a 普通相机或桌面浏览器不能打开小程序：提示在微信中扫码，不改变账号或订单。\\
字段与规则 & AppID、码版本、场景、启停状态与生成时刻；场景只允许公开目录引用。入口展示不证明顾客身份。\\
质量与约束 & 假码展示次数为零；正式扫码成功率须在取得真实平台条件后检验。当前本机仅验证“未就绪时关闭”。\\
\end{BSRUsecase}

\Needspace{11\baselineskip}
\subsubsection{3.2.17 UC-17 验证微信身份并绑定顾客}\label{uc:17}
\begin{BSRUsecase}{UC-17 验证微信身份并绑定顾客}
主执行者 & 顾客。\\
辅执行者 & 微信平台交换临时登录码；BSR 现有顾客账号提供独立密码证明。\\
涉众利益 & 顾客跨端只见本人订单；门店防止错绑、冒领和员工身份混用；管理员可追溯绑定与撤销。\\
触发与前置 & 顾客在小程序发起登录；微信临时码有效，服务端启用真实平台适配；待绑定时顾客知道现有账号凭证。\\
成功后置 & 一个有效微信身份与一个有效顾客账号显式绑定；发行独立于网页 Cookie 的小程序会话，原网页会话不被转移。\\
失败后置 & 码、挑战或账号证明无效时不新建绑定和会话，不泄露目标账号是否有订单。\\
基本流程 & 1 小程序取得一次性临时码并发送服务端。2 服务端向微信平台交换身份，验证 AppID 与响应。3 已绑定身份须复核顾客账号状态及认证纪元后签发小程序会话。4 未绑定身份显示待绑定提示；顾客输入现有账号证明并明确同意。5 服务端核验一次性挑战、账号角色及活跃唯一性，写入绑定和审计，签发会话。\\
扩展流程 & 2a 临时码重放、超时或平台响应不可信：拒绝。3a 顾客停用、须改密或绑定已撤销：旧令牌失效。5a 尝试绑定员工、同一微信身份绑定另一顾客或同一顾客重复绑定：拒绝；换绑须先完成独立撤销与再验证。\\
字段与规则 & AppID、OpenID 的密钥化摘要、内部顾客 ID、一次性挑战摘要、会话令牌摘要、绑定/撤销时刻、认证纪元。原始 OpenID、临时码、会话密钥和令牌不落业务库或日志。\\
质量与约束 & 失败限速、有效期、一次性消耗及对象授权由服务端裁决；正式 AppID 已核对，真实微信登录仍待私有服务端配置、合法请求域名和真机联调。\\
\end{BSRUsecase}

\Needspace{11\baselineskip}
\subsubsection{3.2.18 UC-18 跨端管理同一购物车}\label{uc:18}
\begin{BSRUsecase}{UC-18 跨端管理同一购物车}
主执行者 & 已绑定顾客；网页顾客使用其独立网页会话。\\
涉众利益 & 顾客在两个客户端继续相同选择；门店看到可信规格和租期，不把暂存数量当已承诺库存。\\
触发与前置 & 顾客选择门店、起租/归还日期、类别、尺码、款式和数量；账号有效且有购物车读取权。\\
成功后置 & 服务端以同一顾客 ID 保存购物车版本与明细；另一个客户端重新读取后显示相同选择。\\
失败后置 & 版本冲突不覆盖另一端更新；无效规格或数量不产生部分写入。\\
基本流程 & 1 两端按 Europe/Oslo 的 09:00 起租、16:00 前归还区间查询剩余量。2 顾客筛选装备并改动购物车。3 客户端带当前版本提交。4 服务端核验归属、规格、数量及版本，一次写入新版本。5 两端按本人的授权接口重新读取明细和版本。\\
扩展流程 & 1a 已过日期或区间无货：提示重新选期。4a 另一端先改动：返回冲突与新版本，顾客选择重新合并，不静默覆盖。5a 离线缓存：仅显示待同步草稿，不能宣称库存已锁定。\\
字段与规则 & customer\_id、store\_id、起租日、归还日、规格、数量、version；购物车不锁单件，最终库存由 UC-19 重验。\\
质量与约束 & 共享数据库事实、本人对象授权；不把小程序本地缓存当服务端订单。当前隔离 HTTP 测试已验证跨端版本冲突。\\
\end{BSRUsecase}

\Needspace{11\baselineskip}
\subsubsection{3.2.19 UC-19 跨端核价、同意并预约}\label{uc:19}
\begin{BSRUsecase}{UC-19 跨端核价、同意并预约}
主执行者 & 顾客。\\
涉众利益 & 顾客在任一客户端得到相同的 NOK 正式金额和可解释政策；门店不超卖最后一件；财务可追溯冻结价与押金。\\
触发与前置 & 已登录顾客提交有版本的购物车；租期有效且价目和演示政策已配置。\\
成功后置 & 只创建一张本人预约、价目/政策快照和有效暂留；两端订单 ID 相同。模拟支付状态有明确标识。\\
失败后置 & 旧报价、旧购物车、政策变化、库存冲突或未同意时不占库存、不创建订单。\\
基本流程 & 1 服务端根据购物车核算日租金、每件参考价值 8\% 押金、预计归还时间、超时及损坏规则，签发短时报价凭据。2 小程序/网页显示 NOK 正式金额，定位币种仅作参考换算。3 顾客阅读并勾选当前政策。4 客户端提交报价 ID、购物车版本、确认版本及幂等引用。5 服务端在事务中复核身份、报价指纹、政策、日期段库存和重复请求。6 原子写入预约、明细、政策快照及暂留并返回同一订单。\\
扩展流程 & 1a 价目缺失：停止结账。4a 未阅读或未同意：阻断提交。5a 报价过期、购物车在另端改变、最后一件已被占用或规则更新：不建单，要求重报重确认。5b 同键安全重试：返回原订单，不延长暂留或重复收费。\\
字段与规则 & quote\_id、cart\_version、cart\_fingerprint、政策版本、租金、押金、expires\_at、idempotency\_key、reservation\_id。关键金额均为整数 NOK；8\% 为待 BSR 核准的演示规则。\\
质量与约束 & 报价只读不锁库存；预约事务保证最后一件至多成功一单。当前本机独立令牌/API 已测，正式平台与真机未测。\\
\end{BSRUsecase}

\Needspace{11\baselineskip}
\subsubsection{3.2.20 UC-20 跨端查看本人订单与具名进度}\label{uc:20}
\begin{BSRUsecase}{UC-20 跨端查看本人订单与具名进度}
主执行者 & 顾客。\\
辅执行者 & 网页端租赁专员、技师和经理按授权更新订单同一组业务事实。\\
涉众利益 & 顾客能看备货、交付、逐件归还、鉴定、押金和当前责任岗位；员工不会因另一客户端显示旧缓存而失去责任。\\
触发与前置 & 顾客打开本人订单；其网页或小程序身份有效。\\
成功后置 & 两端按同一订单 ID 读取已提交状态、具名责任与时间；读取本身不修改订单。\\
失败后置 & 他人订单访问被拒绝；断线时旧缓存标过期，不宣称资金已到账。\\
基本流程 & 1 顾客打开本人订单。2 服务端按内部顾客 ID 核对象归属。3 返回冻结的 NOK 金额、取货码、合同/演示支付、备货负责人、取货、逐件归还及鉴定状态。4 网页员工完成具名备货、核码、归还和技师初评/经理复核。5 小程序重新读取共享订单并显示最新责任阶段、可公开的鉴定摘要和照片数量。6 无争议且逐件验收后显示模拟退押过程；真实到账须外部回执。\\
扩展流程 & 2a 未绑定或错绑：拒绝。4a 受损需 24 小时内初评，经理未核费前只显示“待复核”，不得确定扣款。4b 严重损坏超押金：显示联系补差状态，禁止自动从其他渠道扣费。5a 离线：提示上次同步时刻并重试。6a 争议或退款失败：保留未结状态。\\
字段与规则 & reservation\_id、阶段、责任人/岗位、stage\_at、逐件实还时刻、damage\_case、deposit\_settlement；顾客只能见本人允许的照片和费用明细。\\
质量与约束 & 员工操作与顾客读取共享服务端订单；本机 API 测试已验证无损和受损/补差终态，小程序照片原图查看及真实两设备联调仍需证据。\\
\end{BSRUsecase}

\Needspace{11\baselineskip}
\subsubsection{3.2.21 UC-21 退出与解除微信绑定}\label{uc:21}
\begin{BSRUsecase}{UC-21 退出与解除微信绑定}
主执行者 & 已绑定顾客。\\
涉众利益 & 顾客能分别撤销当前小程序会话或主动解除微信关联；其他账号的订单不受影响。\\
触发与前置 & 顾客选择退出或解绑，且当前会话尚有效；解绑另须验证账号证明。\\
成功后置 & 退出只撤销当前小程序令牌；解绑撤销绑定下全部小程序会话并审计，网页独立会话不被冒用。\\
失败后置 & 无效证明或旧会话不能解除他人绑定；已撤销令牌不能重放。\\
基本流程 & 1 顾客选择退出或解绑。2 系统按独立令牌定位当前绑定。3 退出立即撤销本令牌；解绑复核顾客证明后撤销活跃绑定与相关令牌。4 小程序清除私有缓存并回到公开页。\\
扩展流程 & 2a 会话已过期：客户端清除本地私有状态。3a 账号已停用：服务端统一拒绝旧令牌，管理员按受控程序处理绑定。\\
字段与规则 & binding\_id、revoked\_at、unbound\_at、auth\_epoch；登录挑战和临时码不可重复使用。\\
质量与约束 & Web 与 Mini 两种会话分权；不得从小程序退出推断网页已退出。当前只在本机假平台测试，正式微信身份待验证。\\
\end{BSRUsecase}
